\documentclass[a4paper,11pt]{article}
\usepackage[utf8]{inputenc}
\usepackage[spanish]{babel}
\usepackage{amsmath}
\usepackage{amssymb}
\usepackage{geometry}
\usepackage{fancyhdr}
\usepackage{graphicx}
\usepackage{siunitx}
\usepackage{xcolor}

\geometry{margin=1.8cm}
\pagestyle{fancy}
\fancyhf{}
\rhead{Examen 2019 - Soluciones Completas}
\lhead{AIST}
\rfoot{\thepage}

\title{\textbf{Examen Diciembre 2019 - Soluciones Completas\\Architectures et ingénierie des systèmes de transmission}}
\author{}
\date{}

\begin{document}

\maketitle

\section*{Parte A: Sistemas de transmisión RF (Globalstar)}

\subsection*{Datos del sistema}
\begin{itemize}
    \item \textbf{Altitud satélites:} $h = 1414$ km (LEO - Low Earth Orbit)
    \item \textbf{Frecuencias:}
    \begin{itemize}
        \item Satélite $\rightarrow$ Usuario (DL): $f_{\text{DL,S}} = 2.49$ GHz (banda S)
        \item Usuario $\rightarrow$ Satélite (UL): $f_{\text{UL,L}} = 1.62$ GHz (banda L)
        \item Satélite $\rightarrow$ Gateway (DL): $f_{\text{DL,C}} = 6.875-7.055$ GHz (banda C)
        \item Gateway $\rightarrow$ Satélite (UL): $f_{\text{UL,C}} = 5.091-5.250$ GHz (banda C)
    \end{itemize}
    \item \textbf{Satélite:} $\text{PIRE}_{\text{sat}} = -1.5$ dBW, $G_{R,\text{sat}} = 31.8$ dBi, $F_{R,\text{sat}} = 2.7$ dB
    \item \textbf{Usuario:} $\text{PIRE}_{\text{user}} = -4$ dBW
    \item \textbf{Gateway:} $d_{\text{gw}} = 5.5$ m, $\eta = 70\%$, $T_T = 140$ K
    \item \textbf{Distancia Satélite-Gateway:} $d = 4480$ km
    \item \textbf{Ancho de banda sub-canal:} $B = 1.23$ MHz
    \item \textbf{Ancho de banda por canal:} 16.5 MHz (contiene 13 sub-canales FDM)
    \item \textbf{Arquitectura:} 16 haces por satélite (8 polarización circular derecha + 8 izquierda)
\end{itemize}

\subsection*{Pregunta 1: Esquema del sistema}

\textbf{Descripción del sistema Globalstar:}

El sistema consta de tres elementos principales:
\begin{enumerate}
    \item \textbf{Usuario (Terminal móvil):} comunica con el satélite en bandas L (UL, 1.62 GHz) y S (DL, 2.49 GHz)
    \item \textbf{Satélite (LEO a 1414 km):} actúa como repetidor transparente (``bent-pipe'')
    \begin{itemize}
        \item Recibe 16 canales en banda L de usuarios (misma frecuencia, 16 haces)
        \item Transpone a banda C (5-7 GHz) en 8 portadoras FDMA por polarización
        \item Transmite a la Gateway terrestre
    \end{itemize}
    \item \textbf{Gateway (Estación terrestre):} antena parabólica de 5.5 m, comunica en banda C
\end{enumerate}

\textbf{Flujo de comunicación Usuario A $\rightarrow$ Usuario B:}
$$\boxed{\text{Usuario A} \xrightarrow[\text{1.62 GHz}]{\text{UL banda L}} \text{Satélite} \xrightarrow[\text{5.091-5.250 GHz}]{\text{UL banda C}} \text{Gateway}}$$
$$\boxed{\text{Gateway} \xrightarrow[\text{6.875-7.055 GHz}]{\text{DL banda C}} \text{Satélite} \xrightarrow[\text{2.49 GHz}]{\text{DL banda S}} \text{Usuario B}}$$

\textbf{Características clave:}
\begin{itemize}
    \item Técnica MBA (Multiple Beam Array): 16 haces reutilizan la misma frecuencia
    \item CDMA: múltiples usuarios comparten un sub-canal de 1.23 MHz
    \item FDMA en banda C: 8 portadoras separadas en frecuencia por polarización
    \item Duración de conexión: $\sim$14 minutos por satélite (órbita LEO)
\end{itemize}

\subsection*{Pregunta 2: Verificación del ancho de banda en banda C}

\textbf{Objetivo:} Verificar que la banda C puede acomodar 8 portadoras (para 16 canales) + 2 canales de telemetría.

\textbf{Banda C Uplink (Gateway $\rightarrow$ Satélite):}
$$\Delta f_{\text{UL}} = 5.250 - 5.091 = 0.159 \text{ GHz} = 159 \text{ MHz}$$

\textbf{Banda C Downlink (Satélite $\rightarrow$ Gateway):}
$$\Delta f_{\text{DL}} = 7.055 - 6.875 = 0.180 \text{ GHz} = 180 \text{ MHz}$$

\textbf{Ancho de banda necesario:}

Cada canal ocupa 16.5 MHz. Necesitamos:
\begin{itemize}
    \item 8 portadoras (cada una corresponde a 2 canales de polarización opuesta)
    \item 2 canales dedicados a telemetría y comando
\end{itemize}

$$\text{Ancho total requerido} = 10 \times 16.5 = 165 \text{ MHz}$$

\textbf{Verificación:}
\begin{itemize}
    \item \textcolor{red}{Uplink: $165$ MHz $>$ $159$ MHz $\rightarrow$ ¡PROBLEMA! Déficit de 6 MHz}
    \item \textcolor{blue}{Downlink: $165$ MHz $<$ $180$ MHz $\rightarrow$ OK (15 MHz de margen)}
\end{itemize}

\textbf{Solución al problema del UL:} En la práctica:
\begin{enumerate}
    \item Los canales de telemetría usan menos de 16.5 MHz (probablemente $\sim$10 MHz cada uno)
    \item Reducción de bandas de guarda entre canales
    \item Filtros más selectivos permiten solapamiento parcial
\end{enumerate}

\subsection*{Pregunta 3: Efectos de la órbita LEO}

\subsubsection*{Efecto a considerar: Desplazamiento Doppler}

Los satélites LEO se mueven a alta velocidad orbital ($v \approx 7.5$ km/s), causando efecto Doppler:

$$\Delta f_{Doppler} = f_0 \cdot \frac{v_{rel}}{c}$$

\textbf{Para banda L (UL, 1.62 GHz):}
$$\Delta f_{max} = 1.62 \times 10^9 \times \frac{7500}{3 \times 10^8} \approx \pm 40.5 \text{ kHz}$$

\textbf{Para banda S (DL, 2.49 GHz):}
$$\Delta f_{max} = 2.49 \times 10^9 \times \frac{7500}{3 \times 10^8} \approx \pm 62.3 \text{ kHz}$$

\textbf{Para banda C (5-7 GHz):}
$$\Delta f_{max} \approx \pm 150-175 \text{ kHz}$$

\textbf{Consecuencias:}
\begin{itemize}
    \item \textbf{Compensación necesaria:} El receptor debe corregir continuamente el Doppler
    \item \textbf{Métodos:} PLL (Phase-Locked Loop), estimación basada en posición orbital (efemerides)
    \item \textbf{En CDMA:} También afecta sincronización de códigos
    \item \textbf{Variación temporal:} Durante los 14 minutos de conexión, el Doppler varía desde +máx a -máx
\end{itemize}

\subsubsection*{Ventajas LEO vs GEO}

\begin{center}
\begin{tabular}{|l|p{5cm}|p{5cm}|}
\hline
\textbf{Aspecto} & \textbf{LEO (1414 km)} & \textbf{GEO (36000 km)} \\
\hline
\textbf{Atenuación trayecto} & $\sim$160 dB @ 1.6 GHz & $\sim$200 dB @ 1.6 GHz \\
 & \textcolor{blue}{40 dB MENOS} & \\
\hline
\textbf{Potencia TX requerida} & BAJA: -4 dBW (0.4 W) & ALTA: +20 dBW (100 W) \\
 & Terminal portátil OK & Terminal grande \\
\hline
\textbf{Retardo propagación} & $\frac{2 \times 1414}{c} \approx 9.4$ ms & $\frac{2 \times 36000}{c} \approx 240$ ms \\
 & \textcolor{blue}{Imperceptible} & \textcolor{red}{Molesto para voz} \\
\hline
\textbf{Latencia RTT} & $\sim$20 ms & $\sim$500 ms \\
 & Conversación fluida & Eco perceptible \\
\hline
\textbf{Tamaño antena} & Omnidireccional pequeña & Directiva grande \\
 & (móvil, orientación libre) & (fija, tracking necesario) \\
\hline
\textbf{Doppler} & ALTO: ±40-60 kHz & NULO (satélite fijo) \\
 & Compensación obligatoria & No requiere compensación \\
\hline
\textbf{Handover} & Necesario cada 14 min & NO necesario \\
 & Complejidad añadida & Conexión permanente \\
\hline
\textbf{Cobertura/satélite} & Limitada ($\sim$5000 km diám.) & Global (1/3 de Tierra) \\
 & Requiere constelación & 3 satélites cubren todo \\
\hline
\textbf{Vida útil} & 5-8 años & 15+ años \\
 & Reemplazos frecuentes & Mayor estabilidad \\
\hline
\end{tabular}
\end{center}

\textbf{Conclusión:} LEO es óptimo para comunicaciones móviles (baja potencia, baja latencia, terminales portátiles). GEO es mejor para servicios fijos de alta capacidad.

\subsection*{Pregunta 4: Polarizaciones circulares}

\subsubsection*{Enlaces Usuario ↔ Satélite (bandas L/S)}

Se usan \textbf{polarizaciones circulares} (RHCP - derecha, LHCP - izquierda) por:

\textbf{1. Rotación de Faraday en la ionosfera}
\begin{itemize}
    \item La ionosfera (capa de plasma, 50-1000 km altitud) causa rotación del plano de polarización lineal
    \item Ángulo de rotación: $\Omega_F = 2.36 \times 10^4 \frac{B_0 N_e f^{-2}}{\cos(\theta)}$ [radianes]
    \item Proporcional a: densidad electrónica $N_e$, inversamente proporcional a $f^2$
    \item A 1.6 GHz: rotación de varios grados (impredecible, varía con hora del día, actividad solar)
    \item \textbf{Polarización circular es inmune:} Una onda circular rotada sigue siendo circular
\end{itemize}

\textbf{2. Orientación arbitraria del terminal móvil}
\begin{itemize}
    \item El usuario sostiene el teléfono en cualquier orientación
    \item Polarización lineal: pérdida de hasta 20-30 dB si antena perpendicular
    \item Polarización circular: recibe/transmite eficientemente independiente de rotación axial
    \item Solo importa el pointing (dirección), no la rotación alrededor del eje
\end{itemize}

\textbf{3. Reutilización de frecuencia}
\begin{itemize}
    \item 16 haces usan LA MISMA frecuencia (1.62 GHz UL, 2.49 GHz DL)
    \item 8 haces RHCP + 8 haces LHCP
    \item Separación por polarización: aislamiento $>$20 dB
    \item Duplica capacidad espectral sin aumentar ancho de banda
\end{itemize}

\textbf{4. Discriminación de trayectos múltiples (multipath)}
\begin{itemize}
    \item Reflexión invierte el sentido de giro: RHCP $\rightarrow$ LHCP
    \item Receptor RHCP rechaza señal reflejada LHCP
    \item Reduce desvanecimiento por multipercurso
\end{itemize}

\subsubsection*{Enlaces Satélite ↔ Gateway (banda C)}

Pueden usar \textbf{polarizaciones lineales} (H y V) porque:

\textbf{1. Antena fija con orientación controlada}
\begin{itemize}
    \item Gateway tiene parabólica de 5.5 m fija en tierra
    \item Orientación y polarización perfectamente alineadas y conocidas
    \item No hay rotación aleatoria como en terminal móvil
\end{itemize}

\textbf{2. Menor efecto Faraday}
\begin{itemize}
    \item A 5-7 GHz, rotación Faraday $\propto 1/f^2$ es 10× menor que a 1.6 GHz
    \item Gateway puede compensar rotación residual ajustando feed de antena
\end{itemize}

\textbf{3. Mayor eficiencia y simplicidad}
\begin{itemize}
    \item Feeds de polarización lineal son más simples y eficientes para antenas grandes
    \item Menor pérdida de inserción que polarizadores circulares
    \item Menor costo
\end{itemize}

\textbf{4. Mayor capacidad con polarización dual}
\begin{itemize}
    \item H y V ortogonales: aislamiento $>$25-30 dB
    \item Permite transmitir 2× canales independientes simultáneamente
    \item Similar a reutilización con polarización circular
\end{itemize}

\textbf{Nota:} En la práctica, algunos sistemas también usan circular en banda C para uniformidad del diseño.

\subsection*{Pregunta 5: Bilan de liaison UL (Usuario $\rightarrow$ Satélite)}

\textbf{Atenuación en espacio libre:}
$$A_{EL} = 32.45 + 20\log_{10}(d_{\text{km}}) + 20\log_{10}(f_{\text{MHz}})$$
$$A_{EL} = 32.45 + 20\log_{10}(1414) + 20\log_{10}(1620)$$
$$A_{EL} = 32.45 + 63.01 + 64.19 = 159.65 \text{ dB}$$

\textbf{Potencia recibida:}
$$C_{\text{sat}} = \text{PIRE}_{\text{user}} + G_{R,\text{sat}} - A_{EL} - M$$
$$C_{\text{sat}} = -4 + 31.8 - 159.65 - 2 = -133.85 \text{ dBW}$$

\textbf{Potencia de ruido:}
$$T_R = T_0(F - 1) = 290 \times (10^{2.7/10} - 1) = 290 \times 0.862 = 250.0 \text{ K}$$
$$T_T = T_R + T_A \approx T_R = 250 \text{ K}$$
$$N = 10\log_{10}(k \cdot T_T \cdot B) = 10\log_{10}(1.38 \times 10^{-23} \times 250 \times 1.23 \times 10^6)$$
$$N = -168.18 \text{ dBW}$$

\textbf{Relación C/N:}
$$\left(\frac{C}{N}\right)_{\text{UL}} = C_{\text{sat}} - N = -133.85 - (-168.18) = 34.33 \text{ dB}$$

\subsection*{Pregunta 6: Bilan de liaison DL (Satélite $\rightarrow$ Gateway)}

\textbf{Frecuencia central:} $f = 6.965$ GHz = 6965 MHz

\textbf{Atenuación en espacio libre:}
$$A_{EL} = 32.45 + 20\log_{10}(4480) + 20\log_{10}(6965)$$
$$A_{EL} = 32.45 + 73.02 + 76.86 = 182.33 \text{ dB}$$

\textbf{Ganancia antena Gateway:}
$$G_{\text{gw}} = 10\log_{10}\left(\eta \left(\frac{\pi d f}{c}\right)^2\right)$$
$$G_{\text{gw}} = 10\log_{10}\left(0.7 \times \left(\frac{\pi \times 5.5 \times 6.965 \times 10^9}{3 \times 10^8}\right)^2\right)$$
$$G_{\text{gw}} = 10\log_{10}(0.7 \times 4423.7) = 54.91 \text{ dBi}$$

\textbf{Potencia recibida:}
$$C_{\text{gw}} = \text{PIRE}_{\text{sat}} + G_{\text{gw}} - A_{EL} - M$$
$$C_{\text{gw}} = -1.5 + 54.91 - 182.33 - 3 = -131.92 \text{ dBW}$$

\textbf{Potencia de ruido:}
$$N = 10\log_{10}(k \cdot T_T \cdot B) = 10\log_{10}(1.38 \times 10^{-23} \times 140 \times 1.23 \times 10^6)$$
$$N = -171.69 \text{ dBW}$$

\textbf{Relación C/N:}
$$\left(\frac{C}{N}\right)_{\text{DL}} = C_{\text{gw}} - N = -131.92 - (-171.69) = 39.77 \text{ dB}$$

\subsection*{Pregunta 7: C/N Global}

El enlace más limitante es el UL con C/N = 34.33 dB.

\textbf{C/N Global:}
$$\frac{1}{(C/N)_{\text{global}}} = \frac{1}{(C/N)_{\text{UL}}} + \frac{1}{(C/N)_{\text{DL}}}$$
$$\frac{1}{(C/N)_{\text{global}}} = \frac{1}{10^{34.33/10}} + \frac{1}{10^{39.77/10}}$$
$$\frac{1}{(C/N)_{\text{global}}} = \frac{1}{2710.4} + \frac{1}{9482.9} = 4.795 \times 10^{-4}$$
$$(C/N)_{\text{global}} = 10\log_{10}(2084.2) = 33.19 \text{ dB}$$

\subsection*{Pregunta 8: Factor de ruido Gateway}

$$T_T = 140 \text{ K} \Rightarrow F = 1 + \frac{T_T}{T_0} = 1 + \frac{140}{290} = 1.483$$
$$F_{\text{dB}} = 10\log_{10}(1.483) = 1.71 \text{ dB}$$

Esto indica un receptor de muy bajo ruido, típico de estaciones terrestres profesionales.

\newpage
\section*{Parte B: Budget de liaison DWDM}

\subsection*{Pregunta 2: Potencia en punto S}

\textbf{TRX100 (DP-QPSK):}
\begin{itemize}
    \item Número máximo de canales: $N_{\text{ch}} = \lfloor 4000/100 \rfloor = 40$ (grilla 100 GHz)
    \item Potencia amplificador: $P_{\text{out}} = 19$ dBm
    \item Pérdida MUX: $L_{\text{MUX}} = 5$ dB
    \item Pérdida jarretière: $L_J = 0.5$ dB
    \item Potencia por canal en S: $P_{ch,S} = 19 - 5 - 0.5 - 10\log_{10}(40) = 7.49$ dBm
    \item Potencia total en S: $P_{tot,S} = 19 - 5 - 0.5 = 13.5$ dBm
\end{itemize}

\textbf{TRX200 (DP-16QAM):}
\begin{itemize}
    \item Número máximo de canales: $N_{\text{ch}} = 40$
    \item Potencia por canal en S: $P_{ch,S} = 7.49$ dBm
    \item Potencia total en S: $P_{tot,S} = 13.5$ dBm
\end{itemize}

\textbf{TRX400 (DP-16QAM):}
\begin{itemize}
    \item Número máximo de canales: $N_{\text{ch}} = \lfloor 3000/50 \rfloor = 60$ (no caben 80, solo 60)
    \item Con grilla 50 GHz, necesitamos MUX de 80 canales
    \item Pérdida MUX: $L_{\text{MUX}} = 6$ dB
    \item Potencia por canal en S: $P_{ch,S} = 19 - 6 - 0.5 - 10\log_{10}(60) = -5.28$ dBm (PROBLEMA!)
    \item Modo saturado: $P_{tot,S} = 19 - 6 - 0.5 = 12.5$ dBm
    \item Potencia real por canal: $P_{ch,S} = 12.5 - 10\log_{10}(60) = -5.28$ dBm
\end{itemize}

\subsection*{Pregunta 3: OSNR en punto S}

$$\text{OSNR}_S = P_{ch} - P_{ASE} - 10\log_{10}\left(\frac{B_{opt}}{12.5 \text{ GHz}}\right)$$

Donde $P_{ASE} = 10\log_{10}(5.1 \times 10^{-9} \times (G-1))$ con $G = 10^{25/10} = 316.23$

$$P_{ASE} = 10\log_{10}(5.1 \times 10^{-9} \times 315.23) = -55.93 \text{ dBm}$$

\textbf{TRX100/200:} $\text{OSNR}_S = 7.49 - (-55.93) = 63.42$ dB

\textbf{TRX400:} $\text{OSNR}_S = -5.28 - (-55.93) = 50.65$ dB

Sí, debemos considerar el ruido ASE del amplificador de emisión.

\subsection*{Pregunta 4: Potencia mínima en punto R}

La potencia mínima debe satisfacer tanto el umbral de potencia como el umbral de OSNR:

\textbf{TRX100:}
$$P_{ch,R} \geq \max(-14, P_{ch,S} - 10\log_{10}(\text{OSNR}_S/\text{OSNR}_{\text{min}}))$$
Con OSNR$_{\text{min}} = 12 + 2 = 14$ dB:
$$P_{ch,R} = -14 \text{ dBm}$$

\textbf{TRX200:}
Con OSNR$_{\text{min}} = 19 + 2 = 21$ dB:
$$P_{ch,R} = -14 \text{ dBm}$$

\textbf{TRX400:}
Con OSNR$_{\text{min}} = 26 + 2 = 28$ dB:
$$P_{ch,R} = \max(-12, -5.28 - (50.65 - 28)) = \max(-12, -27.93) = -12 \text{ dBm}$$

\subsection*{Pregunta 5: Budget y alcance}

\textbf{Budget:} $B = P_{ch,S} - P_{ch,R} - L_J$

\textbf{TRX100:} $B = 7.49 - (-14) - 0.5 = 20.99$ dB $\Rightarrow$ Alcance = $20.99/0.18 = 116.6$ km

\textbf{TRX200:} $B = 7.49 - (-14) - 0.5 = 20.99$ dB $\Rightarrow$ Alcance = $116.6$ km

\textbf{TRX400:} $B = -5.28 - (-12) - 0.5 = 6.22$ dB $\Rightarrow$ Alcance = $6.22/0.18 = 34.6$ km

\vspace{1cm}
\noindent\rule{\textwidth}{0.4pt}

\begin{center}
\textit{Fin de las soluciones - Examen 2019}
\end{center}

\end{document}
